\documentclass[notes=none]{beamer}

\usepackage{tikz}
\usepackage[linguistics]{forest}
\usepackage[many]{tcolorbox}
\usepackage{comment}
\usepackage{multicol}
\usepackage{listings}

\newif\ifframeinlbf
\frameinlbftrue
\makeatletter
\newcommand\listofframes{\@starttoc{lbf}}
\makeatother

\addtobeamertemplate{frametitle}{}{%
    \ifframeinlbf%
         \addtocontents{lbf}{%
            \protect\makebox[2em][l]{%
                \protect\usebeamercolor[fg]{structure}\insertframenumber\hfill%
            }%
            \protect\hyperlink{page.\insertpagenumber}\insertframetitle\par%
        }%
   \else\fi%
}

\AtBeginSection[]{
  \begin{frame}
  \vfill
  \centering
  \begin{beamercolorbox}[sep=8pt,center,shadow=true,rounded=true]{title}
    \usebeamerfont{title}\insertsectionhead\par%
  \end{beamercolorbox}
  \vfill
  \end{frame}
}

\newcommand{\myalert}[2]{%
    \tikz\node[coordinate](t#1){};\alert<#1>{#2}
}


\newcommand{\myexplain}[3]
{\begin{tcolorbox}
    [   tikznode boxed title,
        enhanced,
        arc=0mm,
        interior style={white},
        attach boxed title to top center= {yshift=-\tcboxedtitleheight/2},
        fonttitle=\bfseries,
        colbacktitle=white,coltitle=red,
        boxed title style={size=small,colframe=red,boxrule=0pt},
        title={\tikz\node[coordinate](s#1) {};#2}]
#3
\end{tcolorbox}
\begin{tikzpicture}[overlay]
\path<#1>[->, thick] (s#1) edge [bend left] (t#1);
\end{tikzpicture}
}

\newcommand{\explain}[2]
{\begin{tcolorbox}
    [   tikznode boxed title,
        enhanced,
        arc=0mm,
        interior style={white},
        attach boxed title to top center= {yshift=-\tcboxedtitleheight/2},
        fonttitle=\bfseries,
        colbacktitle=white,coltitle=red,
        boxed title style={size=small,colframe=red,boxrule=0pt},
        title={#1}]
#2
\end{tcolorbox}
}

\def\point{\begin{tikzpicture}
\fill[orange] (0,0) circle (2pt);
\end{tikzpicture}}

\setbeamertemplate{navigation symbols}{} %Remove navigation bar

\tikzstyle{every picture}+=[remember picture]

\usepackage{color}
\newcommand{\vet}[1]{\foreach \num in {#1}{\el{\num}}}
\newcommand{\el}[1]{\fbox{\parbox[c][2ex][c]{1.1em}{#1}}}
\newcommand{\id}[1]{\parbox[r][2ex][b]{1.7em}{\raggedleft #1}}
\newcommand{\bl}{\color{blue}}
\newcommand{\gr}{\color{gray}}
\newcommand{\re}{\color{red}}

\title{Languages}
\author{Lunjin Lu}
\date{}


\begin{document}
\begin{comment}
\end{comment}
\frameinlbffalse
\begin{frame}[fragile,t,allowframebreaks]{List of Frames}
    \listofframes
\end{frame}
\frameinlbftrue


\frame{ \titlepage
}

\begin{frame}{Recursive Definitions}
A recusive definition defines a set as follows.
\begin{itemize}
  \item Specify basic objects of the set.
  \item Specify how more complex objects
        are constructed from less complex objects. 
  \item State that only objects in the set are the 
  	basic objects and objects  constructed using the specified rules. 
\end{itemize}

\end{frame}

\begin{frame}{Recursive Definition of EVEN}
EVEN is recusively (inductively) defined by the following rules. 
\begin{itemize}
  \item [(1)]  $2\in EVEN$.
  \item [(2)] If $n\in EVEN$ then $n+2\in EVEN$.
  \item [(3)] EVEN doesn't contain any element other than those obtained using rules (1) and (2).
\end{itemize}

The third rule is called the rule of limitation. It is usually omitted 
when a definition is said to be recursive or inductive. 

\end{frame}

\begin{frame}{Membership of EVEN}
Rules in an inductive definition of a set can be used to 
prove that an element is in the set. For instance, 
the following is a proof that $6\in EVEN$.
\begin{itemize}
  \item [(a)] From rule (1), we know that  $2\in EVEN$.
  \item [(b)] From (a) and rule (2), we know that $4\in EVEN$. 
  \item [(c)] from (b) and rule (2), we conclude that $6\in EVEN$. 
\end{itemize}

The rule (2) is used twice in the above reasoning. 
In the first time, $n=2$ and in the second time $n=4$.
\end{frame}

\begin{frame}{Notation}

\end{frame}
\end{document}